\section{The limit of the primary parity obstruction}
\label{sec:complementary-detector-limit}

The ramification condition in Theorem~\ref{thm:odd-ramification}
is an exact boundary for its primary parity detectors.  Beyond that
condition all these detectors vanish, in every matrix rank.  This
does not decide whether the induction ring is polynomial there.

\begin{theorem}\label{thm:complementary-detector-limit}
Let $O$ be a mixed-characteristic discrete valuation ring with odd
residue characteristic $p$, residue field $k$, and normalized valuation
$v$.  Put $e=v(p)$ and suppose
\[
 v_2(e)>v_2(p-1).
\]
Let $E$ be an $O[\epsilon]/(\epsilon^2)$-free $C_p$-lattice,
let $M=E/\epsilon E$, and let $\delta$ be the coefficient connecting
operator on $\widehat H^\bullet(C_p,M)$.  Set $u=\delta^2$.
For every monic irreducible $f\in k[z]$ other than $z$, the
$f$-primary part of $\widehat H^0(C_p,M)$ has even composition
length over $O[u]$.  Thus every nonzero-primary length-parity
detector of Section~\ref{sec:odd-ramification} is zero.

There is no completeness, finiteness, or perfection assumption on
$k$.  The conclusion is a vanishing theorem for these detectors;
it asserts neither generation nor that the reduction kernel is zero.
\end{theorem}

The proof uses the hypersurface operator and periodicity theorem of
\citet[Proposition~1.3, Corollary~1.4, pp.~40--41;
Theorem~6.1(iii), pp.~52--53]{Eisenbud1980}.  We first prove the
geometric parity calculation to which that theorem will be applied.

\subsection{Symmetric factorizations of every rank}

\begin{lemma}\label{lem:symmetric-factorization-parity}
Assume that $O$ is complete and satisfies the hypotheses of the
theorem.  Put $B=O[[x]]$ and $q(x)=(1+x)^p-1$.
If a square matrix $D$ over $B$ satisfies $D^2=q(x)I$, then
\[
 H(D(0))=\ker D(0)/\operatorname{im}D(0)
\]
has finite $O$-length, and that length is even.
\end{lemma}

\begin{proof}
Consider the double cover
\[
 S=B[y]/(y^2-q(x)).
\]
It is a complete local Cohen--Macaulay domain of dimension two.
Indeed it is finite free over $B$, and $q$ has valuation one at
the prime $(x)$ of $B$, so it is not a square in its fraction field.
A maximal Cohen--Macaulay $S$-module is free over $B$: the parameters
$(\pi,x)$ have depth two on it, and the Auslander--Buchsbaum formula
over the regular local ring $B$ gives projective dimension zero.
Conversely a $B$-free $S$-module is maximal Cohen--Macaulay.
Multiplication by $y$ is therefore a matrix of the displayed kind.

After passing to $K=\operatorname{Frac}(O)$ and completing at $x=0$,
the cover becomes $K[[y]]$, because $q'(0)=p$ is a unit in $K$.
The module is torsion-free there and hence free.  Its reduction
modulo $x$ is free over $K[y]/(y^2)$, so $D(0)$ is exact over $K$.
Its homology over $O$ consequently has finite length.

For a maximal Cohen--Macaulay module $V$ define
\[
 \chi(V)=\ell_O H(y:V/xV\longrightarrow V/xV)\pmod2.
\]
An exact sequence of such modules splits over $B$, so reduction
modulo $x$ remains exact.  The resulting one-periodic homology
sequence is an exact three-cycle of finite-length modules.  Its
length defect is twice an image length.  Thus $\chi$ is additive
modulo two.  It is zero on $S$-free modules: modulo $x$, multiplication
by $y$ on $S$ is a contractible length-two nilpotent block.

This homomorphism extends to $G_0(S)$.  Every finite $S$-module has
a second syzygy that is maximal Cohen--Macaulay, by the depth lemma.
Define its value by the alternating sum in a two-step free
presentation.  Schanuel's lemma makes the value independent of the
presentation, since free summands have value zero; a horseshoe
presentation proves additivity.  This argument does not claim finite
projective dimension for arbitrary $S$-modules.

The special fibre is the domain
\[
 S/(\pi)=k[[x,y]]/(y^2-x^p).
\]
Multiplication by $y$ is injective on it, with quotient
$k[[x]]/(x^p)$.  Hence $p[k]=0$ in $G_0(S)$.  Since $p$ is odd,
$\chi(k)=0$, and every finite-length $S$-module has value zero.
The sequence given by multiplication by $\pi$ also gives
$\chi(S/(\pi))=0$.

Let $P$ be a height-one prime of $S$ not containing $\pi$.
For a nonzero element $a(x)+yb(x)$ of $P$, its norm
$a(x)^2-q(x)b(x)^2$ is a nonzero element of $P\cap B$.
Remove its maximal power of $\pi$.  The resulting primitive series
is a unit times a nonconstant distinguished monic polynomial in
$x$.  One can see this preparation directly: modulo $\pi$, write
the series as $x^d\overline U$.  Given a monic degree-$d$ polynomial
$P_d$ and a unit $U$ correct modulo $\pi^j$, divide the error by
$\overline U$ and split it into a polynomial of degree less than
$d$ and $x^d$ times a series.  These two terms correct $P_d$ and
$U$ modulo $\pi^{j+1}$.  Completeness gives the factorization.
The degree is positive because the series lies in a proper prime.
Consequently $S/P$ is finite over $O$.

Its fraction field $L$ is a finite extension of $K$.  The normalization
$O_L$ is finite over $O$, and $O_L/(S/P)$ has finite length.
Write $\alpha,\beta$ for the images of $x,y$ in $O_L$; both lie
in its maximal ideal and $\beta^2=q(\alpha)$.
Over
\[
 S_L=O_L[[x]][y]/(y^2-q(x)),
\]
the point ideal $I=(x-\alpha,y-\beta)$ has quotient $O_L$.
It is free over $O_L[[x]]$ with the two indicated generators,
on which multiplication by $y$ is
\[
 D_{\alpha,\beta}=
 \begin{pmatrix}
 \beta&c(x)\\ x-\alpha&-\beta
 \end{pmatrix},\qquad
 c(x)=\frac{q(x)-\beta^2}{x-\alpha}.
\]
The division is integral since $q(\alpha)=\beta^2$.
This matrix squares to $q(x)I$.  Since $S_L$ is free over $S$,
the point sequence gives $\chi(O_L)=-\chi(I)$.

We compute the specialized length over $O_L$.  Write
$e_L=v_L(p)$, where $v_L$ is normalized.  It is an integer multiple
of $e$, so $v_2(e_L)>v_2(p-1)$.  Suppose first that
$\alpha\beta\ne0$, and put $l=v_L(\alpha)$ and $s=v_L(\beta)$.
Away from the equality threshold $e_L=(p-1)l$, the binomial expansion
gives
\[
 2s=\min(e_L+l,pl).
\]
In the lower range $2s=pl$, so $l$ is even.  In the upper range
$2s=e_L+l$, and $e_L$ is even, so again $l$ is even.
At the threshold, $l=e_L/(p-1)$ is an even integer.

The specialized matrix has entries
$\beta,\beta^2/\alpha,-\alpha,-\beta$.
It is a nonzero rank-one square-zero matrix.  Its homology length
is its nonzero elementary-divisor valuation,
\[
 m=\min(s,l,2s-l).
\]
In the lower range $m=l$; in the upper range $m=\min(e_L,l)$.
At the threshold cancellation can only increase $2s$ above $pl$,
and $m=l$.  All these integers are even.
If $\alpha=0$, then $\beta=0$ and the nonzero entry is $p$,
giving length $e_L$.  If $\beta=0$ but $\alpha\ne0$, then
$1+\alpha$ is a nontrivial $p$th root of unity, and the length is
$v_L(\alpha)=e_L/(p-1)$, also even.  These cases exhaust the points.

Restriction from $O_L$ to $O$ multiplies a length by the residue
degree, so $\chi(I)=0$.  The finite normalization cokernel then gives
$\chi(S/P)=0$.  Finally, prime filtrations generate $G_0(S)$.
The zero prime contributes $S$, the only height-one prime containing
$\pi$ is $(\pi)$, every other height-one prime was treated above,
and the maximal ideal contributes $k$.  All values are zero.
Thus $\chi$ is identically zero, proving the lemma.
\end{proof}

\subsection{From a primary lattice to a symmetric factorization}

\begin{proof}[Proof of Theorem~\ref{thm:complementary-detector-limit}]
We may first complete $O$.  Tate cohomology commutes with this flat
extension, and its finite primary composition series has unchanged
length and residue factors.  Completion also preserves $v(p)$.

Let $T$ be the action on $E$ and lift it to $O[[\epsilon]]$.
Its stable curvature is
\[
 \widetilde U=(\widetilde T^p-I)/\epsilon^2.
\]
To identify its coefficient action, expand the two-periodic Tate
differential as $d(\epsilon)=d+\epsilon Q+\epsilon^2R+\cdots$.
The identity $d(\epsilon)^2=\epsilon^2\widetilde U$ gives
\[
 \widetilde U(0)=Q^2+dR+Rd.
\]
Since $Q$ induces $\delta$, the curvature induces exactly
$u=\delta^2$ on $\widehat H^\bullet(C_p,M)$.
Integral primary decomposition of the curvature therefore selects
the required cohomology pieces.  On every selected nonzero component
the curvature is a unit, and the cancellation argument of
Lemma~\ref{lem:truncation-primary-acyclic}, valid over any complete
DVR, gives zero Tate cohomology of that component of $E$ itself.

We must keep track of primary composition length, rather than only
ordinary $O$-length.  Write the irreducible residue polynomial as
\[
 f(z)=g(z^{p^a}),
\]
with $g$ separable irreducible.  Pass to a finite separable residue
extension splitting $g$, realized by a finite unramified extension
$O'/O$.  Each resulting irreducible factor is
$z^{p^a}-\alpha_j$, of degree $p^a$: a separable extension is
linearly disjoint from the corresponding purely inseparable field
extension.  Each simple original $O[u]$-factor becomes one
$p^a$-dimensional residue factor in each selected summand.  Hence,
on any one such summand $H_j$,
\[
 \ell_{O'}(H_j)=p^a\ell_{O[u]}(H_f).
\]
The factor $p^a$ is odd.  It is enough to prove even ordinary length
on each selected lattice.  Rename $O'$ as $O$ for the next steps;
its value of $e$ is unchanged.

Let $F$ denote a selected curvature-unit lattice.  Since
$\widehat H^\bullet(C_p,F)=0$, the coefficient sequence for
multiplication by the uniformizer gives zero Tate cohomology of
$F/\pi F$.  A finite $k[C_p]=k[x]/(x^p)$-module is a sum of
$k[x]/(x^j)$, $1\le j\le p$; every summand with $j<p$ has
nonzero Tate cohomology.  Thus $F/\pi F$ is free over $k[C_p]$.
Lift a basis.  Nakayama gives a surjection $O[C_p]^r\to F$, and
equality of $O$-ranks gives injectivity.  Hence $F$ is free over
\[
 R=O[C_p]=O[[x]]/(q(x)),\qquad q(x)=(1+x)^p-1.
\]
Multiplication by $\epsilon$ is an $R$-linear matrix $D$ with
$D^2=0$ and $\ker D=\operatorname{im}D$, because $F$ is free
over $O[\epsilon]/(\epsilon^2)$.

If $D$ has a unit entry, choose $v$ with $w=Dv$ primitive.
The reductions of $v,w$ are independent, since $D^2=0$.
Choose a functional $\lambda$ with $\lambda(w)=1$ and
$\lambda(v)=0$.  The projection
\[
 P(z)=\lambda(Dz)v+\lambda(z)w
\]
commutes with $D$ and is the identity on the span of $v,w$.
It splits off the contractible block $Dv=w$, $Dw=0$.
As an $O[\epsilon]/(\epsilon^2)[C_p]$-module this block is free,
so it contributes zero to the coefficient Tate groups.
After removing all such blocks, the remaining exact one-periodic
free complex is minimal over $R$.

Lift the same matrix $D$ in every degree to $B=O[[x]]$ and write
\[
 \widetilde D^2=q(x)V.
\]
Cancellation of $q$ shows that $V$ commutes with $\widetilde D$.
The matrix $V$ is the degree-two hypersurface operator of the minimal
periodic resolution.  By Eisenbud's Theorem~6.1(iii), it gives the
periodicity and is invertible modulo the maximal ideal.  This also
holds for the chosen lift: the operator is unique up to homotopy,
and on a minimal complex a null-homotopic map is zero modulo the
maximal ideal.  Comparison with a matrix-factorization resolution,
where the operator is the identity, proves invertibility of its
residue.  Thus $V$ is a unit matrix over $B$.

After a finite unramified extension, $V$ has a commuting square root.
Here is the residue-field argument without a perfection assumption.
In the finite commutative algebra $k[\overline V]$, each reduced
field factor is $L=k(v)$, with $v\ne0$.  Write its minimal polynomial
as $h(z^{p^b})$ with $h$ separable.  Then
$L_s=k(v^{p^b})$ is separable over $k$.  Adjoin a square root
$\mu$ of $v^{p^b}$ over $L_s$; this is a separable extension.
A finite separable extension of $k$ can contain the required
conjugate copies of these fields and square roots.  In each resulting
factor,
\[
 w=v^{(p^b+1)/2}/\mu\qquad\text{satisfies}\qquad w^2=v.
\]
The root lifts through the nilpotent radical since $2w$ is a unit.
Thus it is a polynomial square root of $\overline V$.
Realize the finite separable residue extension unramifiedly, lift
the polynomial root to $W_0(V)$, and iterate
\[
 W_{j+1}=(W_j+VW_j^{-1})/2.
\]
The matrices converge in the complete local topology to a unit
$W$ with $W^2=V$.  Every iterate, and therefore the limit, commutes
with $\widetilde D$.  Consequently
\[
 D_*:=\widetilde D W^{-1}\qquad\text{satisfies}\qquad D_*^2=qI.
\]
At $x=0$ its kernel and image are those of $\widetilde D(0)$,
because $W(0)$ is a commuting unit.  Unramified extension preserves
ordinary module length.  Lemma~\ref{lem:symmetric-factorization-parity}
therefore makes $\ell_O H(D(0))$ even.

Finally, let $M_F=F/\epsilon F$.  The coefficient sequence
$0\to M_F\to F\to M_F\to0$ identifies $M_F$ with
$\ker D=\operatorname{im}D$.  On the invariants of the free
$R$-module $F$, multiplication by $\epsilon$ is $D(0)$.
Taking invariants gives
\[
 H^1(C_p,M_F)=\ker D(0)/\operatorname{im}D(0).
\]
Since $F$ has zero Tate cohomology, its coefficient connecting map
identifies the two Tate parities of $M_F$.  Their common ordinary
length is therefore even.  Dividing the earlier length identity by
the odd integer $p^a$ proves the required primary parity.
Completion preserves it, so the conclusion holds over the original
DVR as well.
\end{proof}

The theorem explains why higher-rank searches with the same parity
coordinates cannot settle the complementary ramification range.
A classification there requires information beyond these coordinates.
